% Gabarit pandoc des corrigés publiés (brut et guidé). Contenu sans marque ni nom de site.
% Compilation : lualatex (latexmk -lualatex).
% Le Markdown reste la source unique : ne pas éditer ce fichier .tex à la main.
\documentclass[11pt,a4paper]{article}

% --- Moteur, polices (Latin Modern) et langue ---------------------------------
\usepackage{iftex}
\ifLuaTeX\else\errmessage{Ce document se compile avec lualatex}\fi
\usepackage{amsmath,amssymb,mathtools}
\usepackage{fontspec}
% mathrm=sym : \mathrm prend ses glyphes dans la police mathématique (ℓ, μ, Ω), comme KaTeX.
\usepackage[math-style=TeX,bold-style=TeX,mathrm=sym]{unicode-math}
% Repli pour les glyphes absents de Latin Modern (indices et exposants Unicode,
% flèches, filets des schémas en texte, symboles) : DejaVu, puis Latin Modern Math
% (ℓ, lettres mathématiques Unicode des légendes).
\directlua{
  luaotfload.add_fallback("repliTexte", {"DejaVuSerif:mode=harf;", "DejaVuSans:mode=harf;", "file:latinmodern-math.otf:mode=harf;"})
  luaotfload.add_fallback("repliMono", {"DejaVuSansMono:mode=harf;", "DejaVuSans:mode=harf;", "file:latinmodern-math.otf:mode=harf;"})
}
\setmainfont{Latin Modern Roman}[RawFeature={fallback=repliTexte}]
\setsansfont{Latin Modern Sans}[RawFeature={fallback=repliTexte}]
\setmonofont{Latin Modern Mono}[RawFeature={fallback=repliMono}]
\setmathfont{Latin Modern Math}
\usepackage[french]{babel}
\frenchsetup{SmallCapsFigTabCaptions=false}

% --- Mise en page --------------------------------------------------------------
\usepackage[a4paper,margin=2.2cm,headheight=14pt,footskip=1.2cm]{geometry}
\usepackage{microtype}
\setlength{\emergencystretch}{4em}
\tolerance=1500
\setlength{\parindent}{0pt}
\setlength{\parskip}{0.55\baselineskip plus 2pt minus 1pt}
\usepackage{xcolor}
% Gris nommés : babel-french rend « ! » actif dans le document.
\definecolor{gris2}{gray}{0.98}
\definecolor{gris3}{gray}{0.97}
\definecolor{gris5}{gray}{0.95}
\definecolor{gris35}{gray}{0.65}
\definecolor{gris40}{gray}{0.60}
\definecolor{gris45}{gray}{0.55}
\definecolor{gris50}{gray}{0.50}
\definecolor{gris55}{gray}{0.45}
\definecolor{gris60}{gray}{0.40}
\definecolor{gris65}{gray}{0.35}
\definecolor{gris70}{gray}{0.30}
\definecolor{gris80}{gray}{0.20}
\definecolor{gris85}{gray}{0.15}
\usepackage{graphicx}
\usepackage{float}
\usepackage{caption}
\captionsetup{font=small,labelformat=empty,justification=centering}
\usepackage{booktabs,longtable,array,multirow,calc}
\usepackage{siunitx}
\sisetup{locale=FR}
\usepackage[version=4]{mhchem}
\usepackage{listings}
\usepackage{fancyhdr}
\usepackage{lastpage}
\usepackage[most]{tcolorbox}
\usepackage{enumitem}
\setlist{nosep,leftmargin=*,topsep=2pt}
\usepackage[hidelinks,unicode]{hyperref}
\usepackage{bookmark}
\hypersetup{pdftitle={CCINP MP-MPI-PC-PSI-TPC-TSI — Anglais 2026},pdfauthor={Rédigé par IA},pdfcreator={LaTeX},pdfsubject={Corrigé guidé}}
\usepackage{xurl}

% --- Pandoc : listes, tableaux, figures ----------------------------------------
\providecommand{\passthrough}[1]{#1}
\providecommand{\tightlist}{\setlength{\itemsep}{0pt}\setlength{\parskip}{0pt}}

\makeatletter
\def\maxwidth{\ifdim\Gin@nat@width>\linewidth\linewidth\else\Gin@nat@width\fi}
\def\maxheight{\ifdim\Gin@nat@height>0.55\textheight 0.55\textheight\else\Gin@nat@height\fi}
\makeatother
\setkeys{Gin}{width=\maxwidth,height=\maxheight,keepaspectratio}
\setcounter{secnumdepth}{-\maxdimen}
\providecommand{\pandocbounded}[1]{#1}
\usepackage{longtable,booktabs,array}

% --- Code (listings) -----------------------------------------------------------
\lstset{
  basicstyle=\ttfamily\small,
  columns=fullflexible,
  keepspaces=true,
  breaklines=true,
  breakatwhitespace=false,
  postbreak=\mbox{\textcolor{gris50}{$\hookrightarrow$}\space},
  frame=single,
  rulecolor=\color{gris35},
  framesep=4pt,
  xleftmargin=5pt,xrightmargin=5pt,
  backgroundcolor=\color{gris3},
  keywordstyle=\bfseries,
  commentstyle=\color{gris65},
  stringstyle=\color{gris80},
  showstringspaces=false,
  tabsize=4,
  extendedchars=true,
  aboveskip=6pt,belowskip=6pt,
}

% --- Macros des corrigés (KaTeX ↔ LaTeX) ---------------------------------------
% Les corrigés sont écrits pour KaTeX : on définit ici ce que KaTeX connaît et
% que LaTeX n'a pas, sans jamais retoucher le Markdown.
\providecommand{\R}{\mathbb{R}}
\providecommand{\N}{\mathbb{N}}
\providecommand{\Z}{\mathbb{Z}}
\providecommand{\Q}{\mathbb{Q}}
\providecommand{\C}{\mathbb{C}}
\providecommand{\dd}{\mathrm{d}}
\providecommand{\llbracket}{\lBrack}
\providecommand{\rrbracket}{\rBrack}
\providecommand{\lBrace}{\lbrace}
\providecommand{\rBrace}{\rbrace}
\providecommand{\Box}{\square}
\providecommand{\Diamond}{\lozenge}
\providecommand{\varnothing}{\emptyset}
\providecommand{\argmax}{\operatorname*{arg\,max}}
\providecommand{\argmin}{\operatorname*{arg\,min}}
\renewcommand{\leq}{\leqslant}
\renewcommand{\geq}{\geqslant}

% Équation hors texte : réduite à la largeur de la ligne si elle déborde.
\newsavebox{\boiteequation}
\newcommand{\equationajustee}[1]{%
  \sbox{\boiteequation}{$\displaystyle #1$}%
  \ifdim\wd\boiteequation>\linewidth
    \[\resizebox{\linewidth}{!}{\usebox{\boiteequation}}\]%
  \else
    \[#1\]%
  \fi}

% --- Encadrés pédagogiques (sobres, lisibles en noir et blanc) -----------------
\tcbset{
  encadre/.style={
    enhanced,breakable,sharp corners,
    colback=white,colframe=gris55,
    boxrule=0.4pt,leftrule=2.2pt,
    left=6pt,right=6pt,top=3pt,bottom=3pt,
    before skip=6pt,after skip=6pt,
    fonttitle=\bfseries\small,coltitle=black,
    attach title to upper,after title={\quad},
    parbox=false,
  }
}
\newtcolorbox{indication}[1]{encadre,colframe=gris45,colback=gris2,title={Indication~#1}}
\newtcolorbox{methode}{encadre,colframe=gris70,title={Méthode}}
\newtcolorbox{piege}{encadre,colframe=gris85,boxrule=0.8pt,leftrule=2.2pt,title={Piège}}
\newtcolorbox{jury}{encadre,colframe=gris60,colback=gris5,title={Jury}}
\newtcolorbox{encadrenote}{encadre,colframe=gris40}

% --- Titres de section ---------------------------------------------------------
\usepackage{titlesec}
\titleformat{\section}{\Large\bfseries}{}{0pt}{}
\titleformat{\subsection}{\large\bfseries}{}{0pt}{}
\titleformat{\subsubsection}{\normalsize\bfseries}{}{0pt}{}
\titlespacing*{\section}{0pt}{14pt plus 4pt minus 2pt}{6pt}
\titlespacing*{\subsection}{0pt}{10pt plus 3pt minus 2pt}{4pt}
\titlespacing*{\subsubsection}{0pt}{8pt plus 2pt minus 2pt}{2pt}

% --- En-tête et pied de page ---------------------------------------------------
\pagestyle{fancy}
\fancyhf{}
\fancyhead[L]{\small\textsc{Corrigé guidé} · CCINP MP-MPI-PC-PSI-TPC-TSI}
\fancyhead[R]{\small Anglais 2026}
\fancyfoot[L]{\footnotesize Corrigé rédigé par IA, non officiel}
\fancyfoot[R]{\footnotesize \thepage\,/\,\pageref*{LastPage}}
\renewcommand{\headrulewidth}{0.4pt}
\renewcommand{\footrulewidth}{0pt}
\fancypagestyle{plain}{\fancyhf{}\fancyfoot[R]{\footnotesize \thepage\,/\,\pageref*{LastPage}}\fancyfoot[L]{\footnotesize Corrigé rédigé par IA, non officiel}\renewcommand{\headrulewidth}{0pt}}

\begin{document}
\thispagestyle{plain}
\begin{center}
  {\small\scshape Corrigé guidé}\par\smallskip
  {\LARGE\bfseries CCINP MP-MPI-PC-PSI-TPC-TSI — Anglais 2026\par}\medskip
  {\large CCINP · MP-MPI-PC-PSI-TPC-TSI · Anglais · 2026\par}
\end{center}
\vspace{2pt}\hrule\vspace{8pt}

\emph{Corrigé rédigé par IA, non officiel. Signalez toute erreur.}

\hypertarget{pruxe9sentation-du-sujet}{%
\section{Présentation du sujet}\label{pruxe9sentation-du-sujet}}

\textbf{Format.} L'épreuve de langue vivante A dure 3 heures ; elle est commune aux filières MP, MPI, PC, PSI, TPC et TSI et mutualisée avec e3a-Polytech. Toute machine est interdite. Un seul exercice : une \textbf{synthèse en anglais} de quatre documents « d'égale importance », en \textbf{400 mots ± 10 \%} (360 à 440), avec un \textbf{titre obligatoire}, compté dans le total. Le candidat indique le nombre total de mots et trace un trait vertical tous les vingt mots. Dans l'introduction, il ne donne que la \textbf{source et la date} de chaque document ; ensuite, il y renvoie par « Doc.1 », « Doc.2 », etc.

\textbf{Les documents.}

\begin{longtable}[]{@{}
  >{\raggedright\arraybackslash}p{(\columnwidth - 8\tabcolsep) * \real{0.0413}}
  >{\raggedright\arraybackslash}p{(\columnwidth - 8\tabcolsep) * \real{0.1446}}
  >{\raggedright\arraybackslash}p{(\columnwidth - 8\tabcolsep) * \real{0.1640}}
  >{\raggedright\arraybackslash}p{(\columnwidth - 8\tabcolsep) * \real{0.2291}}
  >{\raggedright\arraybackslash}p{(\columnwidth - 8\tabcolsep) * \real{0.3909}}@{}}
\toprule\noalign{}
\begin{minipage}[b]{\linewidth}\raggedright
Doc.
\end{minipage} & \begin{minipage}[b]{\linewidth}\raggedright
Source, date
\end{minipage} & \begin{minipage}[b]{\linewidth}\raggedright
Nature
\end{minipage} & \begin{minipage}[b]{\linewidth}\raggedright
Idée directrice
\end{minipage} & \begin{minipage}[b]{\linewidth}\raggedright
Données utiles
\end{minipage} \\
\midrule\noalign{}
\endhead
\bottomrule\noalign{}
\endlastfoot
1 & thetimes.com, août 2023 & article d'opinion & l'université est devenue une décision financière, à la rentabilité incertaine, surtout pour les femmes et les plus modestes & plus de 60 000 £ de dette (90 000 £ en médecine) ; 28 800 £ pour un diplômé à cinq ans, contre 25 896 £ en moyenne chez les 22-29 ans ; plus de 20 ans pour amortir ; femmes : 3 600 £ (11,8 \%) de moins ; anciens boursiers de cantine : 10 \% de moins ; effacement de la dette repoussé de 30 à 40 ans ; 52 \% rembourseront tout ; frais de 1 000 £ (1998) à 9 000 £ \\
2 & Higher Education Policy Institute, février 2025 & témoignage d'un diplômé & malgré le coût, l'université vaut la peine pour les expériences qu'elle offre & 14,5 heures de travail salarié par semaine ; 27 \% rembourseront leur prêt, 65 \% selon la réforme ; record d'étudiants défavorisés ; autonomie, diversité, stages, compétences relationnelles \\
3 & theguardian.com, août 2024 & enquête, témoignages & la dette pèse sur la carrière, le logement, la santé mentale et les projets familiaux & prêts jusqu'à 120 000 £ ; promotions refusées ; temps partiel au salaire minimum ; diplômés « trapped » et « misled » ; conditions du prêt mal comprises ; une médecin de 43 ans doit encore 36 000 £ \\
4 & Policy Institute, King's College London, mai 2024 & deux graphiques d'opinion (849 diplômés en 2024, 1 222 en 2020) & les diplômés jugent encore leur diplôme utile, moins pour la carrière que pour le savoir et l'expérience & « worth it » : 81 \% pour l'ensemble, 84 \% pour les connaissances (87 \% en 2020), 78 \% pour l'expérience étudiante, 66 \% pour la carrière ; s'ils recommençaient : 48 \% même diplôme, 39 \% autre diplôme, 8 \% pas d'université \\
\end{longtable}

\textbf{Thème et difficulté.} Le corpus porte sur l'enseignement supérieur \textbf{au Royaume-Uni}. Il oppose un coût et une dette croissants (Doc.1, Doc.3) à une valeur personnelle et sociale qui demeure (Doc.2, Doc.4). Les documents sont accessibles ; la difficulté est de les \textbf{croiser} au lieu de les résumer l'un après l'autre, et de lire correctement le Doc.4, qui ne dit pas ce que l'on attend : les diplômés ne regrettent guère l'université, ils regrettent plutôt leur filière.

\textbf{Ce qu'il ne faut pas laisser passer.} Mentionner le Royaume-Uni dans le titre et dans la problématique. Exploiter les quatre documents, y compris les graphiques. Rester neutre : aucune opinion, aucune connaissance extérieure, aucune accroche. Respecter la fourchette de mots, indiquer le total, tracer les traits.

\textbf{Gestion du temps.} Proposition : 40 minutes de lecture et de dépouillement, 25 minutes pour la problématique, le titre et le plan, 80 minutes de rédaction au brouillon puis au propre, 20 minutes pour compter, placer les traits et relire, et une marge de 15 minutes.

\textbf{Ce que dit le rapport 2026.} Le rapport d'anglais (p.~6-14) résume le thème commun : l'université reste « une expérience riche et formatrice », mais son coût croissant rend son bénéfice économique « beaucoup plus incertain » (p.~6). Il juge l'épreuve « clairement formulée et abordable » mais révélatrice de « fortes disparités » (p.~7). Les points \textbf{discriminants} :

\begin{itemize}
\tightlist
\item
  le titre et la problématique, « la partie la plus fragile de l'épreuve » (p.~9), en particulier l'oubli du Royaume-Uni, sanctionné ;
\item
  le croisement, qui « reste la difficulté majeure » (p.~10) ;
\item
  la lecture du Doc.4, objet d'un contresens récurrent (p.~11) ;
\item
  la langue, « critère discriminant » (p.~8).
\end{itemize}

Le rapport ne publie aucune moyenne.

\hypertarget{partie-i-la-synthuxe8se-uxe9tape-par-uxe9tape}{%
\section{Partie I --- La synthèse, étape par étape}\label{partie-i-la-synthuxe8se-uxe9tape-par-uxe9tape}}

Le sujet ne comporte qu'un exercice ; les étapes ci-dessous sont celles du travail, chacune avec sa solution, et la dernière donne la synthèse complète.

\hypertarget{uxe9tape-1-duxe9pouiller-les-quatre-documents}{%
\subsection{Étape 1 --- Dépouiller les quatre documents}\label{uxe9tape-1-duxe9pouiller-les-quatre-documents}}

\begin{indication}{1}

Lisez chaque document en notant à part ses \textbf{idées} (ce qu'il affirme) et ses \textbf{exemples} (chiffres, témoignages). Une synthèse se construit avec les idées ; les exemples ne servent qu'à les appuyer.

\end{indication}

\begin{indication}{2}

Pour le Doc.4, lisez les deux graphiques séparément. Que mesure le premier (le diplôme a-t-il « valu la peine », et pour quelle raison) ? Que mesure le second (que referait-on) ? Comparez 2024 à 2020.

\end{indication}

\begin{indication}{3}

Construisez un tableau dont les lignes sont des \textbf{thèmes} (coût, rentabilité, remboursement, effets sur la vie, accès, valeur de l'expérience) et les colonnes les documents : les cases pleines sur une même ligne sont vos croisements.

\end{indication}

\hypertarget{solution}{%
\subsubsection{Solution}\label{solution}}

Grille de confrontation :

\begin{longtable}[]{@{}
  >{\raggedright\arraybackslash}p{(\columnwidth - 8\tabcolsep) * \real{0.1118}}
  >{\raggedright\arraybackslash}p{(\columnwidth - 8\tabcolsep) * \real{0.1952}}
  >{\raggedright\arraybackslash}p{(\columnwidth - 8\tabcolsep) * \real{0.1893}}
  >{\raggedright\arraybackslash}p{(\columnwidth - 8\tabcolsep) * \real{0.2092}}
  >{\raggedright\arraybackslash}p{(\columnwidth - 8\tabcolsep) * \real{0.2645}}@{}}
\toprule\noalign{}
\begin{minipage}[b]{\linewidth}\raggedright
Thème
\end{minipage} & \begin{minipage}[b]{\linewidth}\raggedright
Doc.1
\end{minipage} & \begin{minipage}[b]{\linewidth}\raggedright
Doc.2
\end{minipage} & \begin{minipage}[b]{\linewidth}\raggedright
Doc.3
\end{minipage} & \begin{minipage}[b]{\linewidth}\raggedright
Doc.4
\end{minipage} \\
\midrule\noalign{}
\endhead
\bottomrule\noalign{}
\endlastfoot
Coût, dette & frais de 1 000 £ à 9 000 £, dette de plus de 60 000 £ & coût de la vie, prêt d'entretien insuffisant, 14,5 h de travail & jusqu'à 120 000 £, gonflés par les intérêts & --- \\
Rentabilité & gain salarial faible ; moindre pour les femmes et les plus modestes & --- & salaires des diplômés qui stagnent & carrière : 66 \% « worth it », 30 \% non \\
Remboursement & effacement à 40 ans au lieu de 30 ; 52 \% rembourseront tout & 27 \% aujourd'hui, 65 \% avec la réforme & conditions mal comprises & --- \\
Effets sur la vie & --- & étudiants pressés par le temps & promotions refusées, logement, santé mentale, famille, sentiment d'être piégé & --- \\
Accès & risque d'une université réservée aux riches & record d'étudiants défavorisés & une médecin conseille aux plus jeunes de s'en passer & --- \\
Valeur de l'expérience & --- & autonomie, résilience, diversité, stages & médecins et juristes ont besoin du diplôme & 81 \% pour l'ensemble, 84 \% savoir, 78 \% expérience ; 8 \% ne retourneraient pas à l'université, 39 \% changeraient de filière \\
\end{longtable}

\begin{methode}

Une ligne remplie dans au moins deux colonnes donne un paragraphe croisé. Une case qui contredit sa voisine (52 \% contre 65 \%, « réservée aux riches » contre « record d'étudiants défavorisés ») est une tension à signaler : c'est ce que le jury appelle « mise en tension ».

\end{methode}

\begin{piege}

Le Doc.1 annonce un effacement de la dette repoussé \textbf{de 30 à 40 ans} (\emph{will now be increased to 40 years}). Le rapport, lui, écrit « de 20 à 30 ans » (p.~10) : c'est le document qui fait foi, et c'est lui qu'il faut restituer.

\end{piege}

\begin{jury}

Rapport 2026 (p.~10) : parmi les points « souvent oubliés ou mal rendus », le jury cite la santé mentale, les inégalités de genre et de milieu social, l'allongement du délai de remboursement et la notion de \emph{tuition fees}, « souvent absente, remplacée par des termes vagues : price, cost, money ». Il reproche aussi d'accumuler les exemples, la vie en colocation par exemple, sans dégager l'idée : « l'université comme espace de transition vers la vie adulte ».

\end{jury}

\hypertarget{uxe9tape-2-probluxe9matique-et-titre}{%
\subsection{Étape 2 --- Problématique et titre}\label{uxe9tape-2-probluxe9matique-et-titre}}

\begin{indication}{1}

Cherchez la tension qui traverse tout le corpus, pas une question factuelle : qu'est-ce qui oppose le Doc.1 et le Doc.3 au Doc.2, et où se place le Doc.4 ?

\end{indication}

\begin{indication}{2}

Formulez la problématique en question \textbf{indirecte} (\emph{They raise the question of whether\ldots{}}) : pas d'inversion, pas d'auxiliaire \emph{do}, pas de point d'interrogation. Mentionnez le Royaume-Uni.

\end{indication}

\begin{indication}{3}

Le titre se choisit \textbf{après} la problématique : un groupe nominal neutre, qui couvre les deux dimensions (coût et valeur) et situe le corpus.

\end{indication}

\hypertarget{solution-1}{%
\subsubsection{Solution}\label{solution-1}}

Problématique : \emph{They raise the question of whether a degree is still worthwhile now that it leaves most graduates heavily in debt.}

Titre : \textbf{The Cost and Value of University Education in the UK} (10 mots). Il nomme les deux pôles du corpus (\emph{Cost}, \emph{Value}), l'objet (\emph{University Education}) et le pays (\emph{the UK}), sans jugement.

\begin{methode}

\emph{worthwhile} (adjectif) évite les fautes classiques sur \emph{worth} : \emph{it is worth it}, \emph{it is worth doing}, mais jamais \emph{it doesn't worth it} ni \emph{education worths}. Le rapport regrette que \emph{worthwhile} soit « rarement utilisé » (p.~13).

\end{methode}

\begin{piege}

Titres refusés par le jury parce que trop vagues, polémiques ou fantaisistes : \emph{University cost}, \emph{Universities in crisis today}, \emph{To study or not to study, that is the question?} (p.~9). Problématiques fautives : \emph{To what extent the university are benefic for student?}, \emph{Should I go in university?}

\end{piege}

\begin{jury}

Rapport 2026 (p.~9) : l'absence de référence au Royaume-Uni « entraîne une pénalisation automatique ». Il faut aussi respecter « l'exigence de neutralité propre à la synthèse », et le rapport demande une « véritable problématique, et non une simple question factuelle » (p.~8).

\end{jury}

\hypertarget{uxe9tape-3-construire-un-plan-croisuxe9}{%
\subsection{Étape 3 --- Construire un plan croisé}\label{uxe9tape-3-construire-un-plan-croisuxe9}}

\begin{indication}{1}

Évitez le plan « inconvénients (Doc.1, Doc.3) / avantages (Doc.2, Doc.4) » : il juxtapose au lieu de croiser. Cherchez une progression.

\end{indication}

\begin{indication}{2}

Une progression possible : de la décision financière (coût, rentabilité, remboursement) à ses effets sur la vie des diplômés, puis à ce que les diplômés retiennent malgré tout de leurs études.

\end{indication}

\begin{indication}{3}

Chaque partie doit citer au moins deux documents, et les quatre documents doivent apparaître dans l'ensemble.

\end{indication}

\hypertarget{solution-2}{%
\subsubsection{Solution}\label{solution-2}}

\begin{longtable}[]{@{}
  >{\raggedright\arraybackslash}p{(\columnwidth - 4\tabcolsep) * \real{0.1359}}
  >{\raggedright\arraybackslash}p{(\columnwidth - 4\tabcolsep) * \real{0.5923}}
  >{\raggedright\arraybackslash}p{(\columnwidth - 4\tabcolsep) * \real{0.2418}}@{}}
\toprule\noalign{}
\begin{minipage}[b]{\linewidth}\raggedright
Partie
\end{minipage} & \begin{minipage}[b]{\linewidth}\raggedright
Idée directrice
\end{minipage} & \begin{minipage}[b]{\linewidth}\raggedright
Documents croisés
\end{minipage} \\
\midrule\noalign{}
\endhead
\bottomrule\noalign{}
\endlastfoot
I & Aller à l'université est devenu une décision financière (frais, dette, coût de la vie) à la rentabilité incertaine & Doc.1, Doc.2, Doc.3 \\
II & Un remboursement plus long et plus lourd façonne la vie des diplômés et menace l'égalité d'accès & Doc.1, Doc.2, Doc.3 \\
III & Les diplômés reconnaissent pourtant à l'université une valeur qui dépasse l'argent & Doc.2, Doc.4, Doc.3 \\
\end{longtable}

Une phrase de conclusion, sans opinion, reprend la tension : une expérience valorisée, un bénéfice financier qui dépend du sexe, du milieu et de la carrière.

\begin{methode}

La troisième partie corrige la lecture hâtive du Doc.4 : seuls 8 \% des diplômés ne retourneraient pas à l'université, tandis que 39 \% choisiraient une autre filière. Le regret porte sur la filière, pas sur l'enseignement supérieur lui-même.

\end{methode}

\begin{piege}

\emph{because of debt, lots of students regret doing studies (Doc.4)} est le contresens que le jury cite (p.~11). Le Doc.4 ne parle d'ailleurs pas de la dette.

\end{piege}

\begin{jury}

Rapport 2026 (p.~10-11) : beaucoup de copies adoptent un schéma descriptif, un document par partie, qui « ne répond pas aux attentes de la synthèse, qui exige de faire dialoguer les documents ». Les meilleures copies « distinguent » une insatisfaction financière d'une remise en cause globale des études (p.~11).

\end{jury}

\hypertarget{uxe9tape-4-ruxe9diger-lintroduction}{%
\subsection{Étape 4 --- Rédiger l'introduction}\label{uxe9tape-4-ruxe9diger-lintroduction}}

\begin{indication}{1}

Trois éléments seulement : la présentation du corpus (source et date, numéros des documents), le thème, la problématique. Pas d'accroche, pas de contexte extérieur.

\end{indication}

\begin{indication}{2}

Gardez l'ordre des documents du sujet et mettez leur numéro entre parenthèses : le lecteur saura ensuite à quoi renvoie « Doc.3 ». Ne citez pas les auteurs.

\end{indication}

\hypertarget{solution-3}{%
\subsubsection{Solution}\label{solution-3}}

\emph{The four documents, from thetimes.com (Doc.1, August 2023), the Higher Education Policy Institute (Doc.2, February 2025), theguardian.com (Doc.3, August 2024) and the Policy Institute at King's College London (Doc.4, May 2024), deal with higher education in the UK. They raise the question of whether a degree is still worthwhile now that it leaves most graduates heavily in debt.}

Soit 58 mots, environ 15 \% du total : avec les 10 mots du titre, il en reste près de 340 pour le développement et la conclusion.

\begin{methode}

\emph{deal with}, \emph{raise the question of whether}, \emph{examine}, \emph{highlight} : quelques verbes d'analyse neutres suffisent. Le rapport recommande \emph{argue}, \emph{highlight}, \emph{point out}, \emph{suggest}, \emph{illustrate} (p.~14).

\end{methode}

\begin{piege}

Les introductions trop longues (parfois plus de 100 mots selon le rapport, p.~9) ou remplies d'éléments extérieurs (\emph{For centuries universities used to be reserved\ldots{}}) sont sanctionnées.

\end{piege}

\begin{jury}

Rapport 2026 (p.~8) : les accroches, acceptables à l'oral, « sont à proscrire à l'écrit ». Et pour les sources (p.~11), il est préférable « de conserver l'ordre donné dans le sujet et de préciser les numéros des documents dès l'introduction ».

\end{jury}

\hypertarget{uxe9tape-5-ruxe9diger-compter-relire}{%
\subsection{Étape 5 --- Rédiger, compter, relire}\label{uxe9tape-5-ruxe9diger-compter-relire}}

\begin{indication}{1}

Ouvrez chaque partie par une phrase qui énonce l'idée, puis appuyez-la sur deux documents, avec un chiffre bien choisi plutôt que cinq.

\end{indication}

\begin{indication}{2}

Reformulez au lieu de recopier : deux ou trois mots cités entre guillemets au plus (\emph{trapped}, \emph{misled}). Variez les connecteurs (\emph{yet}, \emph{while}, \emph{nevertheless}, \emph{accordingly}, \emph{above all}).

\end{indication}

\begin{indication}{3}

Comptez au brouillon, puis placez un trait vertical tous les vingt mots sur la copie, titre compris ; indiquez le total à la fin.

\end{indication}

\hypertarget{solution-4}{%
\subsubsection{Solution}\label{solution-4}}

\textbf{The Cost and Value of University Education in the UK}

The four documents, from thetimes.com (Doc.1, August 2023), the Higher \textbar{} Education Policy Institute (Doc.2, February 2025), theguardian.com (Doc.3, August 2024) and the Policy Institute at King's College London (Doc.4, May \textbar{} 2024), deal with higher education in the UK. They raise the question of whether a degree is still worthwhile now \textbar{} that it leaves most graduates heavily in debt.

Going to university has first become a financial decision. Since 1998, tuition \textbar{} fees have risen from £1,000 to £9,000, so students are likely to graduate with over £60,000 of debt (Doc.1), sometimes \textbar{} up to £120,000 once interest has been added (Doc.3). As the maintenance loan no longer covers the cost of living, \textbar{} many students also work about 14.5 hours a week (Doc.2). Yet the return on this investment is uncertain: five years \textbar{} after graduating, the typical graduate earns only slightly more than the average young worker, while women and students from poorer \textbar{} backgrounds earn even less (Doc.1).

Repayment is also lasting longer, since debts will now be written off after 40 years \textbar{} instead of 30, and between 52\% (Doc.1) and 65\% (Doc.2) of graduates are expected to repay their loans in full. \textbar{} This burden shapes graduates' lives: some turn down promotions or settle for part-time jobs to avoid higher repayments, while the \textbar{} debt makes buying a home harder and weigh on their mental health and family plans (Doc.3). Graduates in lower-paid fields \textbar{} feel ``trapped'' and ``misled'', and most admit they never fully understood the terms of their loans (Doc.3). Doc.1 even fears \textbar{} that university may become reserved for the wealthy, although record numbers of disadvantaged students are still enrolling (Doc.2).

Nevertheless, graduates \textbar{} still value what university brings beyond money. It fosters independence, resilience and people skills, exposes students to people from diverse \textbar{} backgrounds and opens doors to work placements (Doc.2). Accordingly, 81\% of graduates consider their degree worth it overall, and 84\% \textbar{} value the knowledge gained, whereas only 66\% say so about their career prospects (Doc.4). Above all, only 8\% would choose \textbar{} not to go to university again, and 39\% would choose a different degree (Doc.4): regret concerns the choice of course \textbar{} rather than higher education itself, even if one doctor now advises her younger relatives to skip it (Doc.3).

Overall, a \textbar{} university education in the UK remains a valued experience, but its financial benefit increasingly depends on gender, social background and \textbar{} career path.

\textbf{(402 words, title included)}

Le décompte compte chaque suite de caractères séparée par des espaces : \emph{£1,000}, \emph{14.5}, \emph{52\%}, \emph{Doc.1} et \emph{part-time} valent chacun un mot ; les traits verticaux ne comptent pas. Avec 402 mots, la synthèse reste dans la fourchette même si l'on compte autrement quelques expressions.

\begin{methode}

Pour présenter un chiffre sans le commenter, on l'inscrit dans une comparaison qui porte l'idée : \emph{81\% of graduates consider their degree worth it overall \ldots, whereas only 66\% say so about their career prospects}. La subordonnée en \emph{whereas} fait le croisement à l'intérieur même du Doc.4.

\end{methode}

\begin{piege}

Fautes relevées par le jury : \emph{dept} pour \emph{debt}, \emph{refund a loan} pour \emph{repay a loan}, \emph{the coast of life} pour \emph{the cost of living}, \emph{economical viewpoint}, \emph{ancient students} pour \emph{former students} ou \emph{graduates} (p.~12) ; \emph{Since the 1970s, student fees rose a lot}, qui exige le \emph{present perfect} (p.~13). D'où \emph{Since 1998, tuition fees have risen}.

\end{piege}

\begin{jury}

Rapport 2026 (p.~8) : il faut « reformuler les données sans les commenter » et « ne jamais extrapoler, ni introduire de références culturelles externes ». La conclusion est « éventuelle, sans ajout personnel ni information extérieure ». Pour la langue, le jury conseille les collocations du sujet : \emph{student debt}, \emph{tuition fees}, \emph{higher education}, \emph{cost of living crisis} (p.~13).

\end{jury}

\hypertarget{bilan}{%
\section{Bilan}\label{bilan}}

\textbf{Méthode de la synthèse, à retenir.}

\begin{itemize}
\tightlist
\item
  Dépouiller en séparant idées et exemples, puis croiser dans une grille thème × document.
\item
  Problématique en question indirecte, neutre, qui situe le corpus (ici le Royaume-Uni) ; titre choisi après elle, nominal et englobant.
\item
  Plan progressif où chaque partie fait dialoguer au moins deux documents ; tous les documents exploités, y compris les graphiques.
\item
  Introduction courte (source, date, numéro, problématique), aucune accroche, aucun ajout extérieur, aucune opinion.
\item
  Compter, tracer un trait tous les vingt mots, indiquer le total.
\end{itemize}

\textbf{Langue à revoir.} \emph{Present perfect} avec \emph{since} ; interrogatives indirectes sans inversion ; \emph{worth} et \emph{worthwhile} ; -s de la troisième personne et accords (\emph{university is}, \emph{many issues are}) ; prépositions (\emph{go to university}, \emph{in February 2025}) ; collocations (\emph{repay a loan}, \emph{the cost of living}, \emph{tuition fees}).

\textbf{Situations transposables.} Lire un graphique d'opinion sans lui faire dire plus qu'il ne dit (regretter sa filière n'est pas regretter ses études). Croiser une enquête chiffrée avec des témoignages. Construire un plan qui va de la cause (le coût) aux effets (la vie des diplômés) puis à la nuance (la valeur reconnue).

\end{document}
